%==============================================================================
\chapter{Contexte et définition du problème}
%==============================================================================

\section{Présentation générale}
Le projet \textbf{Zümm} vise à concevoir un système d'information
dédié à la gestion et au suivi des activités apicoles. L'apiculture moderne
repose sur la surveillance régulière de nombreuses ruches, souvent réparties
sur plusieurs sites géographiques et exploitées par différents apiculteurs.
Le suivi de la santé des colonies, de leur productivité et des interventions
de maintenance constitue une charge organisationnelle importante, aujourd'hui
gérée de façon largement manuelle et dispersée.

\section{Problématique}
En l'absence d'un outil centralisé, les exploitants apicoles rencontrent
plusieurs difficultés :
\begin{itemize}
  \item absence de traçabilité structurée des visites et des inspections des ruches
  \item difficulté à planifier et à coordonner le travail des agents apiculteurs
  \item manque de visibilité sur la production (miel, poids) et sur la rentabilité de chaque site ou ruche
  \item détection tardive des anomalies (baisse d'effectif, chute de production, ouverture non planifiée d'une ruche)
  \item impossibilité d'exploiter automatiquement les mesures issues de capteurs hétérogènes.
\end{itemize}

\section{Finalité}
Le système doit apporter une réponse structurée en deux volets complémentaires :
\begin{enumerate}
  \item une \textbf{gestion manuelle} des opérations de manutention, de leur
        planification et de leur suivi (première partie, objet du présent projet)
  \item une \textbf{supervision distante et automatisée} des indicateurs de
        performance, alimentée par des capteurs, destinée à être développée dans
        le cadre d'un projet ultérieur mais dont le système doit anticiper
        l'intégration.
\end{enumerate}

